During this project, we have tested the utility of the two
programming languages Scratch and Godot Orchestrator.
Due to the time
constraints we have multiple liability threats that make drawing
decisivie conclusions difficult.
Nevertheless, the results we did get
can still be interpreted.

\subsection{Language Analysis}
Scratch was chosen as it still appears to be the state of the art
programming language for teaching young learners the basics of
programming.
The language is optimized for being user friendly and
comes with a fast browser-based game engine that makes it possible
for the learner to quickly see the resultsof their code and iterate
on it.
The block-based structure is proven to be intuitive and
children learn to use it much faster than they do traditional
programming languages like Java.
However, past studies show that the
children can learn bad habits and get wrong ideas about algorithms,
as they get used to the limited options of Scratch.
Writing code in
block-based languages has also been shown to be slower than in
text-based ones, which for that reason are preferred for implementing
larger programs.
Students may be encouraged to start writing code in
Scratch, but will then transition to a text-based language to learn
more complicated concepts, indicating that Scratch is not suitable
for this.

The flow-based programming paradigm is commonly used in industry by
game designers, audio designers, composors and animators.
In fact,
three of the most popular commercial game engines (Godot, Unity, and
Unreal Engine 5) come pre-installed with a flow-based editor for one
or more of these purposes.
Thus, we assumed that flow-based editors
would be better for complex structures and would be more suitable for
teaching architechture.
For writing code, the most advanced language
we identified was Blueprints for Unreal Engine 5.
This flow based
language is integrated with the engine to the point that even game
programmers will be expected to use tit to implement in-game
behavior.

Both visual programming languages have the same drawback of visual
clutter.
It quickly became clear while using them that large programs
take up a lot of screen space.
This can be worked around by
encapsulating large structures into functions that only take up one
block or node of space.
But it is still a fact that the boxes just
take up more space than a line of code in a text based language.
So
code becomes difficult to read much faster than in text-based
languages.

\subsection{Overall workshop structure}
The workshop was centered around the idea of being a game designer
for a weekend.
The students were tasked with creating a game in which
the implementation of the software architechture concepts we saught
to teach were essential to finishing the game in time.
The workshop
started with a short presentation of the concepts after which the
students were tasked with creating the game.
Due to the low amount of
participants, we had ample opportunity to do one+on+one coaching with
the learners, which led to the recovery fo a lot of qualitative data.

Broadly, we judge the structure to have been succesful.
After the
presentatiton, all the students knew their tasks, and from asking the
students to come up with examples of how the concepts might be
implemented it was clear that the presentation sered its purpose of
giving a simple introduction of the concepts.
The students were
highly motivated to create a tangible product, and became
increasingly independent over time.

However, the concepts we presented were too far beyond the abilities
of the students to grasp them fully before the workshop ended.
Both
of the students using Godot only had time to implement functions and
variables.
While the child using Scratch got way further, none of the
students managed to finish the game in time.
We determined this was
due to us assuming the students had received a more comprehensive
introduction to programming prior to the workshop than turned out to
be the case.

We selected for
children who had already been part of programming workshops in the
past.
We had assumed this meant that they would have a basic
understanding of simple datatypes, loops and if-statements, and how
to understand simple algorithms as a set of instructions that is read
through sequentially.
This turned out to not be the case, which meant
that these introductory concepts had to be introduced before teaching
more complex ones.
This ended up taking a lot of time, and resulted
in the children using Godot not having time to implement
instantiation, complex dataflow, and state management.

\subsection{Learning achievements of the children}

\subsubsection{Quantitative results}

\subsubsection{Qualitative results}

\paragraph{Comparison of the languages used in the workshop}
It was clear that the student using Scratch managed to develop much
more of the game than those using Godot Orchestrator.
Due to the low
amount of participants, it was not possible to clearly determine if
this was due to advantages built into the language.
It seems clear
though that having had prior experience with the language was a huge
advantage.
The students using Orchestrator had to be taught how to
navigate the engine and use the node based language prior to writing
code, while the Scratcher could simply start implementing game
mechanics immediately.
Navigating the user interface of Godot
Orchestrator was a big reason for the game taking so long to
implement, so knowing where every important aspect of Scratch was
located meant that significantly more time could be spent on actually
developing the game.

In spite of the small amount of participants, certain key differences
between Scratch and Godot could still be observed during the
workshop.
These are outlined in detail in the results section.
Broadly speaking, Scratch included multiple abstractions that made it
much easier to use for the young learners who had not yet learned the
necessary mathemathics to compute functionality commonly used in
games.
For example, vectors are commonly used to calculate movement,
and indeed to make the character move using Orchestrator the
programmer is supposed to use vectors.
However, the concept of
vectors is, in Denmark, not taught during the public school years.
As
such, we had to give the learners using Orchestrator an introductory
explanation of the concept before they could use it.
In contrast,
Scratch encodes movement using “move x steps in direction X and y
steps in direction Y”. Children aged 10 and above can generally be
expected to have basic knowledge of movement in coordinate systems,
making the implementation of movement much easier.
Datatypes were
similarly abstracted as Scratch does not let the user define their
own, having them instead b predefined in the blocks the programmer
uses.

A similar issue became clear when it came to navigating the UI and
implementing complex functions.
Scratch has much less functionjality
than the Godot game engine, and as such there are fewer options to
pick from when trying to find the block that will do what you want
when compare to Godot.
Since we had designed the game to be possible
to implement in either engine, this resulted in being an advantage
for the Scratch user who had less options to look through thus taking
less time to implement core features.

While the simplicity of Scratch may be great for letting beginners
quickly create very simple programs, for teaching purposes it is not
without drawbacks.
With no way for the developer to set variable
types, no way for one object to access the variables of another, no
way to change the collission box for in-game objects, no way to do
object pooling, and how the language generally hides complex concepts
behind a layer of abstraction, Scratch does not afford teaching these
concepts.
Orchestrator has all the features of Gdscript, the main
text-based language of the Godot game engine, and is thus much closer
in terms of features to what a profesional developer would use, so it
does not share these problems.
In fact, while designing the game that
the students were to develop, Scratch was continously the limiting
factor in terms of which mechanics could be implemented.

\subsection{Recommendations}

\subsubsection{General recommendations}
Our expectations of the programming knowledge of the children was too high, which
resulted in the workshop being too ambitious.
We recommend that
future researchers and educators take this into account when
designing workshops and classes.
The assumption should be that the
children are novices regardless of having participated in other
coding programs before, and the workshop should account for the time
it takes to teach the basic concepts before teaching advanced ones.
This could be done by teaching over a longer period of time, for
example once a week for a few months.
The first lessons would
encompass basic programming knowledge (loops, if statements, reading
your own code, and understanding an algorithm as a sequence of
instructions for the computer to execute).
Only once the students
have this basic understanding should they be exposed to the more
advanced concepts we were trying to teach.

\subsubsection{Recommendations for researchers}
Repeat the study using more students.

Study whether familiarity with language is more important than simplicity.

Study whether Scratch is still useful when it comes to teaching concepts that the
language abstracts away.

The two main issues with Orchestrator were the UI being difficult to
navigate, and the language being too complex for the students.
Here,
we note that Orchestrator was not our initial best candidate for
testing a node-based language.
It is a community-developed language
for Godot and is not officially supported by the engine or used much
in industry.
Blueprints, the node based language used in Unreal
Engine 5, is an integral part of Unreal Engine 5 and required for use
in development.
It stands to reason that it would be much more
polished and user friendly then, than Orchestrator is.
Unfortunately,
for this project we did not have hardware that supported the
specifications of that Engine and had to settle for the closest
parallel we could find.
Future studies should use Blueprints as the
state of the art node-based code editor for game programming if
possible.

\subsubsection{Recommendations for teachers}
Judging by how
motivated the students were to create the game, and the fact that
they we saw clear signs of them learning the concepts after we had
taught them the prerequisites, we judge this method of teaching as
being succesful, and recommend that it be used to teach young
learners in the future.
Thus we recommend a structured approach to
teach young learners programming concepts as they develop a tangible
product they can relate to.
This workshop sowed that implementing the
functionality of a video game was highly motivating for the children,
and that they were highly motivated to learn and implement the
concepts we were teaching.
Even within the confines of the workshop,
the hildren still found room to experiment, and each game would have
likely ended up looking slightly different had they managed to
finish, even though the studewnts were working with the same set of
concrete instructions.

It seems like familiarity with the used language potentially matters much more than
how the language is constructed.
Learning new concepts as well as a
new language was a large burden that significantly slowed down the
learners using Godot Orchestrator.
Thus, the best language to use for
instructors can not simply be determined by its features.
It is
possible that familiarity matters much more for teaching purposes,
though this would have to be verified in future studies.